\section{Introduction and Scope}

\begin{frame}{Problem Setting}
\begin{itemize}
    \item Closed-loop validation depends on discovering failure-prone scenarios.
    \item In WOMD logs, safety-critical interactions are sparse.
    \item Passive mining is therefore compute-inefficient.
    \item Thesis question: can planner uncertainty be optimized for failure discovery?
\end{itemize}
\end{frame}

\begin{frame}{Why This Matters}
\begin{itemize}
    \item Existing pipelines emphasize collision/risk outcomes, not epistemic difficulty.
    \item If uncertainty objectives are robust, they could improve blind-spot discovery efficiency.
    \item If they collapse, benchmark claims should pivot to more actionable objectives.
\end{itemize}
\end{frame}

\begin{frame}{Scope and Thesis Positioning}
\begin{itemize}
    \item Simulation stack: WOMD/Waymax with WOSAC-aligned evaluation assumptions.
    \item Compared objective families:
    \begin{itemize}
        \item surprise-based counterfactual objectives
        \item risk-based decision objectives
    \end{itemize}
    \item Contribution style: negative-result-aware, reproducibility-first, benchmark-grounded.
\end{itemize}
\end{frame}
